\scanpage{11}
\deleted{Noter d'autre part qu'on a équivalence entre les conditions
suivantes:

a) L'homomorphisme $\phi$ est surjectif (resp. injectif) pour tout $d$
et tout point $s\in U_d$ (définissant un $Y$ correspondant).

b) Même condition, mais avec \uline{un} $d$ fixé multiple de $p$, et tel
que \deleted{\unsure{si}} l'on ait $\operatorname{Ker}v=0$ (condition sans doute automatiquement \unsure{vérifiée}
(8.5.7)).

c) $(\operatorname{Tors}^{n-1}(X)\otimes_W k)_{\mathrm{SS}}=0$ (resp. $(\operatorname{Tors}^{\unsure{n+1}}(X)\otimes_W k)_{\mathrm{SS}}=0$).}

\inserted{8.5.\unsure{11}.} Il serait intéressant de donner des exemples où \deleted{ces conditions ne sont pas
satisfaites.}\inserted{$\phi$ n'est pas bijectif}
\deleted{Noter que la condition non respée est toujours satisfaite
si $n=2$ ($X$ une surface). Dans le cas général,} on a l'inégalité
\[
\inserted{(8.5.1\unsure{0|1}.1)}\quad
\dim_k H^{n-1}(X,\underline O)_{\mathrm{SS}}=\dim_{\underline F_p}H^{n-1}(\bar X,\underline Z/p\underline Z)\geq\dim((H^{\inserted{n-1}}(X)/\operatorname{Tors}^{n-1}(X))\otimes k)_{\mathrm{SS}}
\]
(grâce à 8.1.2 et 8.4.3 ; en fait cela marche en toute dimension $i$, pas
seulement $i=n-1$), et dire que $\phi$ est bijectif signifie que l'inégalité
précédente est une égalité. L'exemple de Serre au symposium de Mexico
fournit un $X$ (quotient d'une intersection complète par un groupe $G=\underline Z/p\underline Z$
opérant sans points fixes) pour lequel on a $H^1(X)=0$, mais $H^1(X,\underline Z/p\underline Z)\simeq\underline Z/$
Quitte à couper avec des hyperplans, on trouve le même exemple avec
$\dim X=2$. Pour cet exemple, \deleted{(8.5.11.1) n'est pas une égalité, et
par suite,} d'après \deleted{ce qui précède,}\inserted{8.1.2. on aura}
\diagram{\begin{tikzcd}[ampersand replacement=\&,column sep=small,row sep=small]
 \& \inserted{H^1(X,\underline O)_{\mathrm{ss}}} \arrow[d,equal] \& \\
({}_{\inserted{p}}\operatorname{Tors}^{\unsure{2}}(X)\deleted{\unsure{\otimes k}})_{\mathrm{SS}} \arrow[r,phantom,"\simeq" description] \& H^1(X,\underline Z/p\underline Z)\otimes k \arrow[r,phantom,"=" description] \& \operatorname{Ker}\varphi
\end{tikzcd}}
(quel que soit $d$ et la section $Y$ de $X$ de degré $d$, permettant de définir
$\phi$), qui est non nul, donc $\phi$ n'est pas injectif. \deleted{Si maintenant $Z$ est
un schéma lisse projectif sur $k$ de dimension $n-2$ ($n$ donné $\geq\unsure{4}$), à cohomo-
logie cristalline sans torsion \inserted{et $(H^{n-4}(Z)\otimes k)_{\mathrm{SS}}\ne0$} (par exemple $P^{n-2}_k$ \inserted{si $n$ pair $\ldots$}), alors, posant
$X'=X\times Z$, on voit (sous réserve de la \uline{formule de Kunneth} en cohomologie
cristalline) que $(\operatorname{Tors}^{n-1}(X')\otimes k)_{\mathrm{SS}}$ \deleted{vient comme facteur direct} est non nul,
ce qui fournit des exemples où $\phi$ n'est pas surjectif, pour tout $n\geq4$.
Il y a sans doute aussi des exemples avec $n=3$, en cherchant bien $\ldots$.}
\mspage{10}
